\documentclass[11pt]{article}

% Anonymous review mode for ARR / EMNLP submission. Switch to "final" for camera-ready.
\usepackage[review]{acl}

\usepackage{times}
\usepackage{latexsym}
\usepackage[T1]{fontenc}
\usepackage[utf8]{inputenc}
\usepackage{microtype}
\usepackage{inconsolata}
\usepackage{graphicx}
\graphicspath{{../../artifacts/figures/}{../figures/}}
\usepackage{amsmath}
\usepackage{amssymb}
\usepackage{enumitem}
\usepackage{booktabs}
\usepackage{xurl}
\usepackage{algorithm}
\usepackage{algpseudocode}
\setlength{\textfloatsep}{8pt plus 2pt minus 2pt}
\setlength{\floatsep}{6pt plus 2pt minus 2pt}
\setlength{\intextsep}{6pt plus 2pt minus 2pt}
\renewcommand\UrlFont{\ttfamily\small}
% Allow URL line breaks at any character (xurl behaviour),
% then provide a non-hyperlink \fpath macro for in-text file paths.
\providecommand{\fpath}[1]{\begingroup\Urlmuskip=0mu plus 1mu\relax\url{#1}\endgroup}

% Override the default "Anonymous ACL submission" placeholder that acl.sty
% emits in [review] mode (see acl.sty L132). "ACL" there refers to the
% Association for Computational Linguistics (the umbrella org), but we make
% the EMNLP 2026 target explicit. Per ACL rules acl.sty itself must not be
% modified --- this is a user-side \renewcommand only.
\renewcommand\outauthor{%
    \begin{tabular}[t]{c}
    \bfseries Anonymous EMNLP 2026 Submission
    \end{tabular}}

\title{EDO-Frame: Auditable Local Organization \\
for Multi-Agent Reasoning}

% Used only in [final] / [preprint] mode; in [review] mode acl.sty replaces
% this with the \outauthor block above.
\author{Anonymous Authors \\
  Submission for EMNLP 2026 Long Paper Track (via ARR) \\
  \texttt{anonymous@anonymous.org}}

\begin{document}
\maketitle
\sloppy

\begin{abstract}
Large Language Model (LLM) multi-agent systems commonly rely on either a central orchestrator or fixed specialist roles. These assumptions make systems easy to control, but they obscure a basic question: how can useful division of labor form when agents have only local visibility, no global expert catalog, and no predeclared identities? We introduce \textbf{EDO-Frame}, an auditable framework for local multi-agent organization. Agents operate on sparse neighborhoods, choose among answering, delegating, splitting, auditing, tool use, and memory-update actions, and update future routing state from accepted downstream value rather than self-reported competence. We deliberately separate this full framework from delivered evidence. The executed Stage-1 system, \textbf{Terminal-Consensus Peer Backpropagation (TCPB)}, implements only fixed topologies, role-prior nodes, linear handoff packets, and terminal outcome calibration. On HotpotQA chain-200 with \texttt{gpt-4.1-mini}, TCPB is not the strongest F1 method: \texttt{self\_claim} reaches 0.7641 F1 while TCPB reaches 0.7381 at the same token budget. A separate local \texttt{Phi-4-mini} diagnostic shows the opposite capability-boundary behavior: \texttt{edo\_stage2\_chain} improves F1 over \texttt{single\_agent} by +7.93 points on the same 200-example slice, with a single-seed paired-bootstrap 95\% CI of [+1.52, +14.39] points, but at 2.33$\times$ token cost. The contribution is therefore a claim-bounded framework, a reproducible negative strong-backbone Stage-1 measurement, and an initial local-model evidence path; multi-dataset, multi-seed, and external-baseline closure remains pending before stronger performance claims are justified.
\end{abstract}

\section{Introduction}

The dominant design patterns in LLM-based multi-agent systems make a strong assumption before work begins: either a central orchestrator decides who should do what \citep{wu2023autogen,hong2024metagpt,shen2023hugginggpt}, or each agent receives a fixed specialist identity in advance \citep{qian2024chatdev,chen2024agentverse}. These shortcuts are controllable, but they leave open a basic organizational question. In realistic collaboration, useful division of labor is rarely known at initialization; it forms through partial knowledge, repeated interaction, local judgment, and uneven structural exposure.

This paper takes that organizational view seriously. We ask whether near-homogeneous language agents, each with only local neighbor visibility on a sparse graph, can form useful task division through recursive delegation, recursive acceptance, and value-based social updating. Our object is not merely more dialogue turns or a stronger router; it is how tasks move through a constrained society and how local trust, outsourcing, splitting, and audit induce recognizable organization.

The failure mode motivating our work remains \textbf{delegation miscalibration}. Prior systems often assume that an agent can accurately assess whether it should accept a task or hand it off. In practice, LLMs are prone to overconfident self-assessment \citep{guo2017calibration,kadavath2022know}: they often accept work they should not keep, delay delegation, or pass forward low-quality intermediate results. Earlier versions of our project showed that replacing self-reflection \citep{shinn2023reflexion,madaan2023selfrefine} with terminal outcome feedback significantly improves delegation safety on fixed-topology benchmarks. However, we now argue that this is only the first layer of a richer theory. Delegation quality is not just a routing score problem; it is a problem of organizational emergence under local interaction.

We therefore propose \textbf{Emergent Delegation Organization (EDO)} as the full long-paper method. In EDO, agents start from the same cognitive machinery and the same action grammar. The only allowed asymmetries arise from graph position, local interaction history, and the task neighborhoods each agent encounters over time. Every task is treated as a node in a task tree, not merely as a single linear message. Upon receiving a task, an agent chooses among three primitive actions: \texttt{do\_self}, \texttt{outsource(neighbor)}, or \texttt{split(subtasks)}. Importantly, each outsourcing decision also creates a recursive accountability chain: if agent $a$ delegates to $b$, then $a$ is responsible for accepting, revising, rerouting, or rejecting $b$'s returned result. This recursive audit structure turns delivered value, not self-description, into the basis of social reputation.

Under this framing, an agent's ``personality'' is not a stylistic prompt trait but a structured operational tag vector induced by interaction history. Future local delegation decisions compare task demands against these socially formed tags under local visibility constraints, asking whether useful division of labor can emerge from local interaction alone. We explicitly separate theory from currently executed evidence: our present codebase implements a restricted Stage-1 system, \textbf{Terminal-Consensus Peer Backpropagation (TCPB)} --- fixed \texttt{chain} and \texttt{star} topologies, role-prior nodes, linear handoff packets, and terminal-only competence updates on the final accepting node. We reinterpret that system as the first executable approximation of the broader organizational framework, not as the final method.

\begin{figure*}[t]
\centering
\includegraphics[width=0.84\textwidth]{figure1.png}
\caption{Conceptual overview. EDO-Frame reframes multi-agent reasoning from single calls, fixed roles, debate, or central planning into auditable local organization with recursive delegation, local audit, governed memory, and tag-gated tools; the diagram is explanatory rather than empirical evidence.}
\label{fig:paradigm-comparison}
\end{figure*}

\textbf{Core contribution.} \textbf{EDO-Frame turns decentralized multi-agent reasoning into a testable organizational problem: agents must form routing, audit, memory, and tool-use structure from local interaction history rather than from a predeclared global role table.}

Our contributions are four-fold:
\begin{enumerate}[leftmargin=*,itemsep=0pt,topsep=2pt]
  \item We reframe decentralized multi-agent routing as a problem of \textbf{organizational emergence under local interaction}, rather than merely better static role assignment or router tuning.
  \item We introduce \textbf{EDO-Frame}, a method built around sparse local graphs, recursive \texttt{do / outsource / split} decisions, governed note-board memory, tag-gated tool access, recursive upstream audits, and value-induced personality tags.
  \item We formalize the distinction between the \textbf{full framework} and the \textbf{current restricted executable instantiation}, positioning TCPB as a Stage-1 prototype rather than the endpoint of the research line.
  \item Using current HotpotQA evidence, we report claim-bounded measurements: TCPB is F1-dominated by \texttt{self\_claim} on a capable API backbone, while a local \texttt{Phi-4-mini} diagnostic shows a single-seed paired-bootstrap F1 gain for \texttt{edo\_stage2\_chain} at higher token cost.
\end{enumerate}

\section{Related Work}

\subsection{Routing and orchestration in multi-agent systems}

Much of the LLM multi-agent literature is organized around central decomposition and routing. \textbf{AutoGen} \citep{wu2023autogen} provides a generic conversation framework where a manager agent dispatches sub-tasks to specialist agents; \textbf{MetaGPT} \citep{hong2024metagpt} hard-codes a software-engineering pipeline (PM $\to$ architect $\to$ engineer) where each agent's role and order are pre-assigned; \textbf{HuggingGPT} \citep{shen2023hugginggpt} uses an LLM controller to plan, select, and execute models from a global tool catalog. These orchestrator-based frameworks achieve strong control and interpretability, but they share three properties our work specifically removes: (i) a globally visible expert/role table at initialization, (ii) a central decision node that owns all routing, and (iii) a fixed task script. \textbf{ChatDev} \citep{qian2024chatdev} and \textbf{AgentVerse} \citep{chen2024agentverse} extend this paradigm with richer communication but still rely on manager-like or globally synchronised state. Recent multi-hop QA systems keep the same explicit-control flavor: \textbf{MA-RAG} \citep{nguyen2025marag} hard-codes a planner/step-definer/extractor/answerer retrieval stack, while \textbf{ReAgent} \citep{zhao2025reagent} adds rollback-based trajectory correction. Our question is narrower and more organizational: how can division of labor emerge when no node has a global expert table and each node sees only its local sparse-graph neighborhood?

\subsection{Agent memory, tools, and protocol boundaries}

Agent engineering frameworks increasingly expose memory stores, tool registries, and structured tool-call protocols as first-class affordances. EDO-Frame uses those ideas only as \emph{governed local state}: each agent receives a small core toolset, chooses specialised tools from a tag-gated pool, writes note-board memory through append-only events, and emits parseable action envelopes. This differs from giving every agent a large global tool stack. In our setting, memory and tools are part of the orchestration mechanism being audited, not hidden external advantages. The current prototype therefore implements a dependency-light local \texttt{NoteBoardMemory}, \texttt{ToolRegistry}, \texttt{ToolSelector}, and action/event schema rather than importing a broad product-agent runtime.

\subsection{Reflection, critique, and outcome-based calibration}

Reflection has become a common mechanism for improving LLM reasoning. \textbf{Reflexion} \citep{shinn2023reflexion} introduces verbal RL where an agent critiques its own trajectory and rewrites a policy. \textbf{Tree-of-Thoughts} \citep{yao2023tot} backtracks across self-evaluated reasoning branches. \textbf{Self-Refine} \citep{madaan2023selfrefine} iteratively rewrites outputs based on the model's own feedback. At the multi-agent level, \textbf{Multi-Agent Debate} \citep{liang2023mad,du2024improving} and \textbf{ChatEval} \citep{chan2024chateval} replace self-critique with peer critique that a judge or meta-reviewer aggregates into the current-turn answer. TCPB/EDO use a different supervision object: no critique transcript is merged into the answer; only accepted-value signals on delegation edges update future routing statistics across episodes. Full EDO keeps this edge-level object via recursive upstream audits in a task tree, so the target of learning is routing state rather than current-answer text; a controlled MAD head-to-head is deferred to \S\ref{sec:experiments} E3.

\subsection{Division of labor, reputation, and local organization}

Our long-paper framing is closer to a theory of organization than to static role play. Classical organization theory \citep{galbraith1973designing,mintzberg1979structuring} establishes that stable roles in human groups are outcomes of repeated coordination, evaluation, and structural position rather than perfect advance assignment, and that reputation is socially constructed through accepted contributions rather than self-claims. Local neighborhood structure also matters: sparse connectivity creates bridge positions, clustered exposure, and bounded trust propagation. EDO imports these intuitions into language-agent systems by making local delegation, recursive acceptance, and value-based personality formation the primary objects of study.

\section{Methodology}
\label{sec:methodology}

\begin{figure}[t]
\centering
\includegraphics[width=0.96\columnwidth]{figure2.png}
\caption{Single-column EDO-Frame pipeline sketch. A task packet becomes local agent views; agents interact through sparse delegation, governed note-board memory, tag-gated tool selection, action envelopes, and audit events. The figure is a mechanism overview, not a reported result.}
\label{fig:edo-pipeline}
\end{figure}

\subsection{Problem formulation}

We model a multi-agent society as a sparse connected directed graph $\mathcal{G} = (\mathcal{A}, \mathcal{E})$, where $\mathcal{A} = \{a_1, a_2, \dots, a_N\}$ is the set of agents and $\mathcal{E}$ defines allowed communication channels. The full EDO framework targets near-homogeneous initialisation where agents share base cognitive machinery and action grammar, with only mild asymmetry from graph position and interaction history. The delivered Stage-1 TCPB prototype is more restricted: it is decentralized within a fixed role-prior topology with hand-tuned safety priors (decomposer, evidence seeker, verifier, synthesizer, and a hand-coded force-forward gate). Near-homogeneous initialisation with role-priors removed is reserved for Stage-2 (\S\ref{sec:experiments} E1).

Each input instance $x = (q, y, C)$ consists of a question $q$, gold answer $y$, and optional supporting context $C$. In the full method, $x$ expands into a \textbf{task tree} rather than a single linear message. Each task node records its parent, children, current uncertainty, required output, evidence state, owner agent, executor agent, candidate result, and audit status. This representation is necessary because a task may be solved directly, outsourced intact, or recursively split into subtasks that later require integration and audit.

\textbf{Decentralization invariant (1-locality).} $\pi_i(z)$ factors through $S_i^{t-1}$ and $\{S_j^{t-1}: j\in\mathcal{N}(i)\}$ only (formal statement and graph-class verification are provided in the supplementary appendix).

\subsection{Agent state and personality}

At time $t$, each agent $a_i$ maintains state
\[
S_i^t = \{P_i^t, B_i^t, M_i^t, H_i^t, \mathcal{N}(i)\},
\]
where $P_i^t$ is the agent's public personality-tag vector, $B_i^t$ is its local belief state over neighbors, $M_i^t$ is its private episodic memory, $H_i^t$ is a summary of past delegation and audit interactions, and $\mathcal{N}(i)$ is its visible neighbor set.

The key object is $P_i^t$, the \textbf{value-induced operational persona}. This is not a stylistic role prompt. It is a structured vector whose components summarize socially observed capability, for example:
\begin{multline*}
P_i = [\text{solve}, \text{decompose}, \text{audit}, \text{integrate},\\
\text{explore}, \text{efficiency}, \text{reliability}].
\end{multline*}
These components are intended to capture how useful the agent has been to others in the society: whether its outputs are accepted, whether it reduces or increases rework, whether its decompositions help, whether its integrated outputs are reliable, and whether it achieves useful work efficiently.

Each agent also maintains only \textbf{local} beliefs about neighbors. There is no global competence directory. For a neighbor $a_j$, the local belief $B_i^t(j)$ is constructed from interaction history with $a_j$, recent publicly visible persona summaries, and local audit outcomes involving $a_j$. This preserves decentralization while allowing trust and specialization to be gradually formed.

\subsection{Task signatures and local utility}

Instead of matching tasks to fixed role names, EDO represents each task node $z$ by a task-signature vector
\[
\begin{aligned}
\phi(z)=[&\text{decompose},\text{verify},\text{integrate},\text{explore},\\
&\text{breadth},\text{uncertainty},\text{cost}].
\end{aligned}
\]
This task signature can be extracted through rules, light-weight LLM judgment, or future learned modules. The key point is conceptual: routing should be based on the match between task demands and socially formed personality tags, not on a hard-coded role lookup and not necessarily on a large embedding router. In the current executable protocol, $\phi(z)$ is extracted deterministically from multi-hop cues, hop depth, evidence count, uncertainty, and hop budget (supplementary appendix).

We define the demand--persona matching function as the clipped normalised inner product $\mathrm{Fit}(P, \phi){=}\max\big(0,\;\langle P,\phi \rangle / (\lVert P\rVert\cdot\lVert\phi\rVert)\big)\in[0,1]$, with the convention $\mathrm{Fit}(0,\cdot){=}\mathrm{Fit}(\cdot,0){=}0$; the clip makes the range well-defined even when $P,\phi$ admit negative components (the raw cosine is in $[-1,1]$, but a negative match is treated as ``no match'', not as a repulsive signal). In Stage-1 TCPB, $P$ and $\phi$ collapse onto a scalar competence $c\in[0,1]$ (mean-axis fold, \S 3.6) and are non-negative by construction, so $\mathrm{Fit}$ degenerates to a one-dimensional comparison and the clip is inactive.

For an agent $a_i$ receiving task node $z$, the three candidate utilities are:
\begin{align*}
U_i^{\mathit{self}}(z) ={} & \mathrm{Fit}(P_i, \phi(z)) \\
                           & - \lambda_c \mathrm{Cost}_{\mathit{self}}(z) \\
                           & - \lambda_r \mathrm{Risk}_{\mathit{self}}(z), \\
U_i^{\mathit{out}}(z, j) ={} & \mathrm{Fit}(B_i(j), \phi(z)) \\
                            & - \lambda_s \mathrm{SendCost}(i,j) \\
                            & - \lambda_a \mathrm{AuditCost}(z) \\
                            & - \lambda_r \mathrm{RejectRisk}(j,z), \\
U_i^{\mathit{split}}(z) ={} & \mathrm{SplitGain}(z) - \lambda_m \mathrm{MergeCost}(z) \\
                            & - \lambda_d \mathrm{DepthPenalty}(z) \\
                            & - \lambda_a \mathrm{AuditLoad}(z).
\end{align*}
The selected action is
\begin{align*}
\pi_i(z) = \arg\max \big\{ & U_i^{\mathit{self}}(z), \;\; U_i^{\mathit{split}}(z), \\
                            & \textstyle\max_{j \in \mathcal{N}(i)} U_i^{\mathit{out}}(z,j) \big\},
\end{align*}
subject to budget, depth, and safety constraints.

\subsection{Recursive upstream audit}

The main departure from the current prototype is that delegation is not terminal at handoff time. Every outsourcing decision also creates an \textbf{acceptance obligation}. If $a$ delegates a task to $b$, then $a$ must later evaluate whether to accept $b$'s result, reject it and redo the task, reject it and reroute, or reject it and resplit.

This produces a recursive structure. If $a \rightarrow b \rightarrow c$, then $b$ audits $c$ before returning upward, and $a$ audits $b$. The system therefore operates on a \textbf{delegation tree with a recursive acceptance ladder}, not merely on a one-dimensional chain of messages. This is essential because social reputation should be grounded in accepted delivered value, not only in terminal answer correctness; the supplementary appendix gives an $\Omega(H)$ information-theoretic credit-assignment argument for the advantage of R2 over terminal-only audit.

\textbf{Complexity bound on Process($z$).} With the Stage-2 bounds $(H_{\max}, d_{\max}, n_{\max}, k_{\max}){=}(4,3,12,3)$, the per-input LLM-call count is bounded by
\begin{align*}
T_{\mathrm{LLM}}(z)
&\le \min\{(H_{\max}{+}1)k_{\max}^{d_{\max}},\\
&\hspace{2.6em} n_{\max}(H_{\max}{+}2)\}=72 .
\end{align*}
In TCPB (Split disabled, Audit$\equiv$Accept), this collapses to $T_{\mathrm{LLM}}^{\mathrm{TCPB}}(z)\le H_{\max}{=}4$, matching the observed chain-200 count (derivation in the supplementary appendix).

\subsection{Personality-tag update}

Each audit event generates a local social signal for the downstream executor. For a parent-upstream node $u$, downstream node $v$, and task node $z$, a local audit event can be written as:
\begin{multline*}
\ell_{u \rightarrow v, z} = [\text{accepted}, \text{rework\_cost},\\
\text{value\_gain}, \text{timeliness},\\
\text{decomposition\_help}, \text{integration\_help}].
\end{multline*}
After the root task finishes, the system also observes a terminal event
\begin{multline*}
r_{\mathrm{terminal}}(x) = [\text{quality}, \text{total\_cost},\\
\text{total\_depth}, \text{final\_accept}].
\end{multline*}
The full update rule combines local accepted-value signals with slower terminal correction:
\begin{align*}
P_i^{t+1} = \mathrm{clip}\Big(&(1-\mu) P_i^t + \mu \big(\eta_{\mathrm{loc}} \mathrm{LocalValue}_i^t \\
&+ \eta_{\mathrm{trm}} \mathrm{TerminalValue}_i^t \\
&- \eta_{\mathrm{rew}} \mathrm{ReworkPenalty}_i^t \big)\Big).
\end{align*}

\textbf{Convergence of persona-tag update (stationarity).} Under (A1) i.i.d.\ stationary signals with $\mathbb{E}[V_i^t]{=}\bar V_i$ and (A2) axis-wise boundedness, the EMA recurrence $P_i^{t+1}{=}(1{-}\mu)P_i^t{+}\mu V_i^t$ converges in $L^2$ to the unique fixed point $P_i^\infty{=}\bar V_i$ (proof in the supplementary appendix; non-stationary regime tracks a sliding window with horizon $\approx 1/\mu$).

\subsection{Governed memory and tool access}
\label{sec:memory-tools}

The governed action-loop design is easy to miss in a pure routing formulation: the agent is not only choosing a neighbor, but also choosing what state and tools are allowed to enter the next hop. EDO-Frame therefore treats memory and tools as auditable routing inputs rather than as an unbounded global capability stack.

The local memory object is a \textbf{note board}: an append-only event log with scope labels (\texttt{task\_public}, \texttt{agent\_private}, \texttt{neighbor\_belief}, and \texttt{tool\_policy}) and mandatory source references. Retrieval is deterministic in the current implementation: tag overlap, entity overlap, lexical overlap, accepted-evidence bonus, and recency tie-breaks determine what a local agent can see. The tool object is a \textbf{tag-gated selector}: core tools are always available, while specialised tools are scored by tag fit, agent prior, observed success, cost, and risk. The selector supports \texttt{tag\_gated}, \texttt{all\_tools}, \texttt{random\_k}, and \texttt{static\_k} modes for future ablations.

These components exist as a dependency-light local prototype under the anonymous code supplement. Unit tests pass locally and remotely; memory writes carry source references, selector outputs log score components, and events serialize to JSONL-compatible dictionaries. A Phi-4 n=50 transparent-bridge diagnostic shows no chain regression but also no F1/EM movement under memory/tool/drift toggles, because those signals are logged rather than injected into planner/checker prompts. Thus this layer remains \emph{method infrastructure}, not component-level performance evidence.

\begin{figure}[t]
\centering
\includegraphics[width=0.96\columnwidth]{figure3.png}
\caption{Governed action-loop sketch for one local handoff. Memory retrieval, tool gating, action envelopes, downstream response, and audit feedback remain conceptual until consumption-layer ablations show non-zero effects.}
\label{fig:action-loop}
\end{figure}

\subsection{Current restricted executable instantiation: TCPB prototype}

The full EDO Stage-2 execution loop is a recursive \textsc{Process} function that selects among the three primitive actions (R1 \textsc{DoSelf} / \textsc{Outsource} / \textsc{Split}), performs per-handoff audit (R2) on every \textsc{Outsource} return, and updates the vector belief $B_i^{t+1}(j)$ via an EMA step (R3) for each audit event in the resulting task tree. The supplementary appendix gives complete pseudocode, including the recursion bounds ($H_{\max}{=}4$, $d_{\max}{=}3$, $n_{\max}{=}12$, $k_{\max}{=}3$), the four-class \textsc{Audit} decision protocol, and the Stage-1 fallback projection. The reported empirical results in Table~\ref{tab:main-results} come from the degenerate Stage-1 instance of this loop, which we call \textbf{Terminal-Consensus Peer Backpropagation (TCPB)}: \textsc{Split} is disabled, \textsc{Audit} defaults to \textsc{Accept} on every return, and the vector belief $B_i^t(j)$ is collapsed onto a scalar $c[i]$ via mean-axis fold. With these collapses, the Stage-1 TCPB outsourcing utility reduces to the explicit weighted form
\[
\begin{aligned}
U^{\text{out}}_i(z, j) ={}& 0.55\,c[j] + 0.20\,\mathrm{acc}(j)\\
&+ 0.10\,\mathrm{fb}(j) - \lambda_a\,\mathrm{audit}(z),
\end{aligned}
\]
where $c[j]\in[0,1]$ is the scalar competence projection of $B_i(j)$, $\mathrm{acc}(j)$ is $j$'s historical acceptance rate as upstream auditor, $\mathrm{fb}(j)$ is a topology-driven forward-bias prior (set to $0$ for the Stage-1 chain topology; non-trivial only for Stage-2 sparse graphs), and $\lambda_a{=}0.02$ is the accept margin. The Prototype Scope Box below restates this restriction component-by-component.

\textbf{Reduction (TCPB as a degenerate EDO).} $\mathcal{E}_{\mathrm{TCPB}}$ is obtained from $\mathcal{E}_{\mathrm{EDO}}$ by (R$_\textsc{Split}$) disabling \textsc{Split} via gate weight $-\infty$, (R$_\textsc{Audit}$) $\textsc{Audit}{\equiv}\textsc{Accept}$, and (R$_B$) projecting the 7-d belief onto the scalar $c[i]$ via mean-axis fold; Table~\ref{tab:main-results} numbers are produced under this reduction (full schema in the supplementary appendix).

\textbf{Stage-1 operational collapses.} In the TCPB reduction, \textsc{Split} is disabled and \textsc{Audit} defaults to \textsc{Accept}, so the Stage-2 split/audit hook terms collapse to zero. The remaining quantities become fixed audit cost $\lambda_a{=}0.02$, zero chain send cost, token-proxy self cost $\tau_z/10^3$, self risk $1{-}c[i]$, scalar competence $c[i]{=}\tfrac{1}{7}\sum_k B_i(j)[k]$, and zero forward bias for the chain topology. Hand-set multipliers and update rates are listed in the supplementary appendix; ties prefer \textsc{DoSelf}, then \textsc{Outsource}, then \textsc{Split}, with equally scored neighbours ordered by index.

\subsection{Why the separation matters}

Without separating the full EDO method from the current TCPB prototype, the paper falls into one of two traps. If we write only what the code already does, the paper becomes a fixed-role routing paper with limited conceptual novelty. If we write only the grand theory and ignore implementation limits, reviewers will rightly reject the paper for overclaiming. The correct scientific stance is to make EDO the main theory and TCPB the current restricted instantiation. The exact boundary between the two is fixed by the Prototype Scope Box that follows.

\medskip
\noindent\fbox{%
\begin{minipage}{0.95\columnwidth}
\small
\textbf{Prototype Scope Box} (single source of truth for what the current submission instantiates).

\smallskip
The current executable system instantiates only a restricted Stage-1 subset of EDO. Concretely, the prototype replaces:
\begin{enumerate}[label=(\roman*),leftmargin=*,topsep=2pt,itemsep=1pt]
  \item arbitrary sparse graphs with fixed \texttt{chain}/\texttt{star} topologies;
  \item near-homogeneous agents with four role-prior nodes (\texttt{decomposer}, \texttt{evidence\_seeker}, \texttt{verifier}, \texttt{synthesizer});
  \item task-tree state with a linear \texttt{HandoffPacket};
  \item the three-action \texttt{do\_self / outsource / split} EDO policy with deterministic \texttt{accept-or-forward};
  \item recursive persona-tag updates with terminal-only scalar competence updates (TCPB).
\end{enumerate}
The deterministic instantiations of $\phi(z)$, the utility estimators $U^{self}, U^{out}$, and the persona update rule for this prototype scope are given in the anonymous executable specification supplement; that note is the canonical companion specification for everything in this box and is referenced by Sections~\ref{sec:methodology}, \ref{sec:experiments}, and \ref{sec:limitations} without restating its content. Stage-2 work will progressively replace each of (i)--(v) with the full EDO mechanism; none of those replacements are claimed by the current submission.
\end{minipage}}
\medskip

\subsection{Stage-2 Roadmap (boundaries the current submission does not cross)}
\label{sec:stage2-roadmap}

The supplementary appendix records the full EDO Stage-2 loop. Stage-2 work activates three mechanisms that the present submission deliberately disables or projects: \textbf{R1} recursive \textsc{Split}, \textbf{R2} four-class upstream \textsc{Audit}, and \textbf{R3} seven-axis vector-belief updates. Each is tied to a falsifiable hypothesis: split should help deeper-hop tasks, audit should reduce terminal error propagation, and vector belief should produce non-random specialization while preserving answer quality. Operational definitions and rejection criteria live in the executable supplement; empirical numbers in Section~\ref{sec:experiments} therefore characterise only the restricted Stage-1 / early local-chain settings, not the full R1/R2/R3 system.

\section{Experiments}
\label{sec:experiments}

\subsection{What current evidence already establishes}

Our current reproducible evidence comes from the Stage-1 TCPB prototype on HotpotQA. The main benchmark slice is a fixed set of 200 HotpotQA validation questions, frozen as a versioned slice in the anonymous code/data supplement. Existing runs compare \texttt{fixed\_peer\_calibrated}, \texttt{fixed\_self\_calibrated}, \texttt{fixed\_self\_claim}, \texttt{fixed\_static\_roles}, and centralized baselines under shared inputs and logged execution traces.

These experiments do \emph{not} establish the full EDO theory. What they do establish is narrower but still important:
\begin{itemize}
  \item decentralized outcome-based calibration matters;
  \item delegation safety is measurable and improvable;
  \item structured intermediate state is necessary for route-level auditing;
  \item route preferences can stabilize without a central controller;
  \item some remaining errors are generated by local answer bottlenecks rather than by routing failure.
\end{itemize}
This evidence should be presented as support for the \textbf{direction} of the theory, not as proof of the entire organizational framework.

\subsection{Stage-1 setup: current prototype experiments}

\begin{itemize}[leftmargin=*,itemsep=0pt,topsep=2pt]
  \item \textbf{Benchmark:} HotpotQA \citep{yang2018hotpotqa} distractor validation subset, 200 examples (versioned slice in the anonymous supplement). MuSiQue \citep{trivedi2022musique} and 2WikiMultiHop \citep{ho20202wikimultihop} are reserved for the Stage-2 EDO benchmark transfer (\S 4.4 E5).
  \item \textbf{Model:} \texttt{gpt-4.1-mini} via an OpenAI-compatible API endpoint (held constant across all runs; integrity-check provenance is captured in \S Limitations item~(5)), temperature 0, chain topology unless otherwise noted.
  \item \textbf{Canonical run:} a single Stage-1 batch covering \texttt{fixed\_peer\_calibrated}, \texttt{fixed\_static\_roles}, and \texttt{fixed\_self\_claim} on the 200-sample slice, all validated by an automated log-completeness checker. Metrics consolidated in the per-run \texttt{main\_table.csv}.
  \item \textbf{Current runtime:} an asynchronous parallel runner (\texttt{runner.py}), a method registry encoding all baselines (\texttt{methods.py}), and a typed inter-agent packet schema (\texttt{contracts.py}); released in the anonymous code supplement.
  \item \textbf{Current method objects:} fixed topology, role-prior nodes, linear handoff packet, deterministic local accept-or-forward policy, terminal-only outcome update.
  \item \textbf{Current process metrics:} EM, F1, mean handoff count, dead-end rate, premature accept rate, forward-after-correction rate, heuristic and API token cost, and cost-normalized F1.
  \item \textbf{Note on prior GLM results:} All Stage-1 runs prior to the backbone upgrade used \texttt{glm-4-flash}. These are archived as Stage-1 guidance-only artifacts (in the anonymous supplement). They established the direction of the research but are not the final paper evidence.
\end{itemize}

\subsection{Current main findings from the restricted instantiation}

The canonical Stage-1 \texttt{gpt-4.1-mini} results ($n=200$, chain) support four main observations, summarised in Table~\ref{tab:main-results}.

\begin{table}[t]
\centering
\small
\setlength{\tabcolsep}{3pt}
\begin{tabular}{@{}lcccr@{}}
\toprule
\textbf{Method} & \textbf{EM} & \textbf{F1} & \textbf{MHC} & \textbf{Tok.} \\
\midrule
\texttt{self\_claim}     & 0.635 & \textbf{0.7641} & 2.00 & 6{,}414 \\
\texttt{static\_roles}   & 0.610 & 0.7454 & 1.68 & \textbf{4{,}663} \\
\texttt{peer\_calibrated}& 0.615 & 0.7381 & 2.00 & 6{,}414 \\
\bottomrule
\end{tabular}
\caption{Stage-1 results on HotpotQA chain-200 with \texttt{gpt-4.1-mini}. MHC = mean handoff count; Tokens = mean API tokens per sample. All three methods achieve premature-accept rate $\mathrm{PAR}=0.0$.}
\label{tab:main-results}
\end{table}

\begin{table}[t]
\centering
\small
\setlength{\tabcolsep}{3pt}
\begin{tabular}{@{}lccr@{}}
\toprule
\textbf{Local method} & \textbf{EM} & \textbf{F1} & \textbf{Tok.} \\
\midrule
\texttt{single\_agent}      & 0.265 & 0.3646 & \textbf{1{,}762} \\
\texttt{edo\_stage2\_chain} & \textbf{0.315} & \textbf{0.4440} & 4{,}106 \\
\bottomrule
\end{tabular}
\caption{No-paid local open-weight diagnostic on the same HotpotQA chain-200 slice with \texttt{Phi-4-mini}. The chain improves F1 by +7.93 points over \texttt{single\_agent}; a 10,000-resample paired bootstrap gives a single-seed 95\% CI of [+1.52, +14.39] F1 points and sign-test $p{=}0.0186$. EM is positive but not significant; token cost is 2.33$\times$ higher.}
\label{tab:local-phi4}
\end{table}

\paragraph{Finding 1: Backbone capability changes the value of organization.} All three Stage-1 methods improve by $+15$ to $+20$ F1 points when switching from \texttt{glm-4-flash} to \texttt{gpt-4.1-mini} on the same 200-sample slice and identical routing code. On the capable API backbone, routing is secondary to generation quality; on the local \texttt{Phi-4-mini} diagnostic, the chain produces a positive F1 gain but pays a large token-cost premium (Table~\ref{tab:local-phi4}).

\paragraph{Finding 2: Method ordering is backbone-sensitive and must not be over-interpreted from a single backbone.} On \texttt{glm-4-flash}, the ordering was \texttt{static\_roles} $>$ \texttt{peer\_calibrated} $\approx$ \texttt{self\_claim}. On \texttt{gpt-4.1-mini}, it inverts: \texttt{self\_claim} (0.7641 F1) $>$ \texttt{static\_roles} (0.7454) $>$ \texttt{peer\_calibrated} (0.7381). On \texttt{Phi-4-mini}, the local Stage-2 chain beats \texttt{single\_agent} in F1 but not in cost-normalized F1. This supports a capability-boundary reading rather than a universal ``more agents is better'' claim.

TCPB should \emph{not} be described as the strongest Stage-1 method on the strong backbone: \texttt{self\_claim} leads F1, while \texttt{static\_roles} is most cost-efficient. Full-validation reruns, external baselines, and multi-seed local replications are still required before a finalized ranking.

\paragraph{Finding 3: PAR collapses to zero regardless of backbone.} All three methods achieve PAR$={=}0.0$ on \texttt{gpt-4.1-mini}, consistent with the GLM result. The PAR$=0$ outcome is attributable to the decomposer force-forward gate (the generator-side safety gate), a fixed safety prior that already saturates PAR for \texttt{self\_claim} and \texttt{static\_roles}; the TCPB peer-calibration mechanism itself does not provide additional PAR reduction in the delivered Stage-1 configuration. The safety \emph{gate} is robust to backbone change; TCPB's PAR contribution on top of the gate is not demonstrated here.

\paragraph{Finding 4: delegation overhead helps only under specific capability and cost regimes.} \texttt{peer\_calibrated} matches \texttt{self\_claim}'s token budget on the strong backbone but with lower F1; in contrast, the local chain improves F1 on \texttt{Phi-4-mini} while more than doubling token use. The supported claim is therefore conditional: EDO explains when organization may compensate for individual model limits, not that delegation is uniformly Pareto-superior.

\noindent\textbf{Evidence still needed.} Best-Paper readiness still requires multi-seed local replication, SmolLM3-3B control, MuSiQue / 2Wiki transfer, causal consumption-layer ablations, a manual error taxonomy, and matched MA-RAG / ReAgent / MAD rows. The transparent-bridge ablation is negative for component-causality claims, so these gaps remain closure tasks rather than completed results.

\section{Conclusion}

We proposed \textbf{EDO}, treating multi-agent language systems as organizations formed via delegation and value-based updating. The current evidence is deliberately bounded: TCPB is F1-dominated by \texttt{self\_claim} on a capable API backbone, while a local \texttt{Phi-4-mini} diagnostic shows a single-seed paired-bootstrap F1 gain for the chain at higher cost. The paper's strongest current contribution is therefore not a universal performance claim, but an auditable framework and evidence path for testing when local organization becomes valuable.

\section*{Limitations}
\label{sec:limitations}

We explicitly acknowledge seven scope limitations of the present submission. (1) The current runtime implements only the restricted Stage-1 TCPB prototype: fixed topologies, role-prior nodes, linear packet passing, deterministic accept-or-forward routing, and terminal-only outcome updates. It therefore cannot yet establish the strongest claims of the full EDO framework --- recursive split value, local audit superiority, or personality-tag emergence from near-homogeneous initialization. (2) Some current routing improvements remain partially entangled with hand-designed safety priors and answer-generation constraints; the present evidence should be interpreted as validating a claim-bounded capability-boundary hypothesis, not as proof that local interaction alone suffices for rich organizational emergence across arbitrary settings. (3) Backbone-sensitive method ordering: on our 200-sample HotpotQA chain slice, the relative ranking of delegation methods differs across \texttt{gpt-4.1-mini}, \texttt{Phi-4-mini}, and earlier weak-backbone diagnostics. We \emph{conjecture} that organizational advantage is most visible on harder tasks, weaker individual agents, or richer topologies where single-agent self-assessment breaks down; this conjecture remains under-tested. (4) Statistical evidence is still incomplete: Table~\ref{tab:local-phi4} includes a single-seed paired-bootstrap CI for the local F1 gain, but the paper still lacks multi-seed variance, multiple datasets, and paired statistics for all headline comparisons. (5) \textbf{Cross-endpoint comparability is an open variable}: a provider-side data-integrity event during this submission cycle led to a deliberate switch of the upstream OpenAI-compatible API channel serving the strong backbone, with affected reruns documented in the supplementary appendix. (6) \textbf{Delivered-system Pareto-domination}: the delivered TCPB Stage-1 instance is Pareto-dominated by \texttt{fixed\_self\_claim} on F1 at equal token cost on the chain-200 HotpotQA slice with \texttt{gpt-4.1-mini}; the positive local-chain result is not Pareto-dominant because it costs 2.33$\times$ more tokens. (7) \textbf{Demographic and societal scope}: the only benchmark used (HotpotQA) is sourced exclusively from English Wikipedia and skews toward Western, biographical, and factual entities; the present evaluation does not test multilingual generalisation, low-resource language behaviour, or culturally diverse query distributions.

These limitations are not accidental flaws in reporting. They define the boundary between the current restricted instantiation and the long-paper method it motivates. The next-stage implementation burden is therefore clear: introduce task-tree state, recursive audit, persona-tag updates, and general sparse-graph support, then test whether the stronger organizational predictions of EDO hold empirically.

\section*{Ethical Considerations}
\label{sec:ethics}

\paragraph{Data.} The sole benchmark used (HotpotQA) is sourced from publicly released Wikipedia article excerpts; the dataset contains no personally identifiable information about article subjects beyond what is already publicly available, and no additional PII was collected, generated, or processed. No human subjects were recruited or evaluated in this work.

\paragraph{Potential misuse.} The Emergent Delegation Organization (EDO) framework describes how to assemble multiple LLM agents into a cooperative multi-hop reasoner; in principle the same scaffolding could be repurposed for adversarial routing (e.g., routing queries through agents with hidden specialisations to bypass single-agent safety filters). We note two mitigating observations: (a) the delivered Stage-1 TCPB instance is Pareto-dominated by a trivial single-agent baseline on the canonical strong backbone (\S Limitations~(6)), which reduces its adversarial utility in its present form; (b) the generator-side safety gate is explicitly a safety prior rather than a routing safety claim, and we report negative ablation evidence for it as the attribution source of PAR$=0$ (\S 4.3 Finding~3), which avoids overclaiming safety as an emergent routing property.

\paragraph{Fairness and scope of claims.} As reported in \S Limitations~(7), the single benchmark used skews toward English Wikipedia and Western, biographical, and factual entities; we explicitly do not claim that the EDO organisational framework, the TCPB Stage-1 calibration mechanism, or the proposed Stage-2 R1/R2/R3 mechanisms generalise to multilingual, low-resource, or culturally diverse query distributions. Broader societal-impact studies (e.g., stratified error analysis along demographic axes of queried entities) are an open direction for Stage-2 evaluation.

\paragraph{Reproducibility and release.} Runtime code (\texttt{runner.py}, \texttt{methods.py}, \texttt{contracts.py}), the exact 200-question seed slice, validation scripts, per-run logs, and the full LLM prompt templates are committed for release under the MIT license in the anonymous code/data supplement cited in the supplementary appendix. We do not release model checkpoints because the system is API-only. The provider-integrity event is disclosed in the supplementary appendix to allow independent replication under an equivalent integrity-guarded channel.

\paragraph{AI-assisted writing and coding.} Per the supplementary appendix, code-completion assistance from a commercial LLM assistant was used during development; all resulting code was reviewed and validated by the authors. Manuscript text was authored by the authors; LLM assistance, when used for phrasing, was restricted to grammar and style passes without introducing new claims or citations.

\bibliography{custom}

\end{document}
